% Source: https://artofproblemsolving.com/community/c3h1958532_2004_amc_10a_problem_21
% Stable identity: aops-topic-1958532-q01
Let $\theta$ be the acute angle. The two acute sectors together make up the fraction $\dfrac{2\theta}{2\pi}=\dfrac{\theta}{\pi}$ of each region, and the obtuse sectors the fraction $1-\dfrac{\theta}{\pi}$.

\begin{asy}
usepackage("amsmath");
size(5cm);
defaultpen(fontsize(10pt));
fill(arc((0,0),3,-12.857142857142858,12.857142857142858)--arc((0,0),2,12.857142857142858,-12.857142857142858)--cycle,gray(0.75));
fill(arc((0,0),3,167.14285714285714,192.85714285714286)--arc((0,0),2,192.85714285714286,167.14285714285714)--cycle,gray(0.75));
fill(arc((0,0),2,12.857142857142858,167.14285714285714)--arc((0,0),1,167.14285714285714,12.857142857142858)--cycle,gray(0.75));
fill(arc((0,0),2,192.85714285714286,347.14285714285717)--arc((0,0),1,347.14285714285717,192.85714285714286)--cycle,gray(0.75));
fill((0,0)--arc((0,0),1,-12.857142857142858,12.857142857142858)--cycle,gray(0.75));
fill((0,0)--arc((0,0),1,180-12.857142857142858,180+12.857142857142858)--cycle,gray(0.75));
draw(circle((0,0),1));draw(circle((0,0),2));draw(circle((0,0),3));
draw(3.3*dir(12.857142857142858)--3.3*dir(192.85714285714286));draw(3.3*dir(-12.857142857142858)--3.3*dir(167.14285714285714));
draw(arc((0,0),2.5,-12.857142857142858,12.857142857142858),Arrows(TeXHead));label("$\theta$",(2.5,0),W,fontsize(9pt));label("$3$",3*dir(60),NE);label("$2$",2*dir(75),N);label("$1$",(0,0.95),S,fontsize(8pt));
label("\mbox{}",(min(currentpicture,user=true)+max(currentpicture,user=true))/2,black); // keeps the text after the figure black when LaTeX typesets the labels inline
\end{asy}

The outer ring has area $9\pi-4\pi=5\pi$, the middle ring $4\pi-\pi=3\pi$ and the inner disc $\pi$. So the shaded area is
\[
5\pi\cdot\dfrac{\theta}{\pi}+3\pi\left(1-\dfrac{\theta}{\pi}\right)+\pi\cdot\dfrac{\theta}{\pi}=3\theta+3\pi .
\]
The total area is $9\pi$ and shaded $=\dfrac{8}{13}\,$unshaded, so shaded $=\dfrac{8}{21}\cdot9\pi=\dfrac{24\pi}{7}$. Then $3\theta+3\pi=\dfrac{24\pi}{7}$ gives $\theta=\boxed{\textbf{(B)}\ \dfrac{\pi}{7}}$.
